% sections/background.tex — Background
\section{Background}
\label{sec:background}

\subsection{The Actor Model}
\label{sec:bg:actor}

Hewitt, Bishop and Steiger~\citep{hewitt1973actors} introduced the
actor model in 1973 as a foundation for describing concurrent
computation in artificial intelligence systems.  The model rests on
three primitives: an actor processes messages from its mailbox one
at a time, may send messages to actors whose identifiers it knows,
and may create new actors.  No shared state, no locks.  Agha's
1986 dissertation~\citep{agha1986actors} gave the first formal
treatment, casting actor systems as labeled transition systems and
establishing the equivalences that justify mailbox semantics.
Subsequent work added typed actors~\citep{he2014typed},
capabilities~\citep{clebsch2015pony}, and supervision
trees~\citep{armstrong2007erlang}; the latter is the closest
ancestor of \system{}'s supervisor design.

\subsection{GPU Programming Conventions}
\label{sec:bg:gpu}

Stub: this subsection will summarise the kernel-launch model
(grid/block/thread hierarchy, host-device memory transfer,
synchronization conventions), the SIMT execution discipline, and
the evolution from CUDA~3.0 (kernels as functions, no persistence)
through CUDA~9.x (cooperative groups) to CUDA~12.x on Hopper
(thread block clusters, DSMEM, async memory pools).

\subsection{Persistent Kernels and Their Limits}
\label{sec:bg:persistent}

Stub: persistent kernel pattern from
\citet{gupta2012persistent,aila2009raytrace,steinberger2014whippletree};
why each prior approach stops short of the actor abstraction
(typically because the kernel is treated as a single producer–consumer
loop rather than a population of independently-supervised actors).
